\chapter{Enquadramento nos Objetivos de Desenvolvimento Sustentável (ODS)}
\label{cap:ods}

A Agenda 2030 para o Desenvolvimento Sustentável, adotada em 2015 pelos 193 Estados-membros da Organização das Nações Unidas (ONU), estabelece 17 Objetivos de Desenvolvimento Sustentável (ODS) e 169 metas voltadas a promover o desenvolvimento econômico, social e ambiental integrado \cite{onu2015agenda2030}. Este capítulo fundamenta o enquadramento do problema de monitoramento e manutenção colaborativa de infraestrutura urbana nos compromissos da Agenda 2030, detalhando o ODS principal, suas metas e indicadores oficiais, bem como os dois ODS secundários aos quais o sistema se articula.

\section{ODS Principal: ODS 11 --- Cidades e Comunidades Sustentáveis}
\label{sec:ods-principal}

O problema abordado neste trabalho enquadra-se de maneira central e direta no \textbf{Objetivo de Desenvolvimento Sustentável 11 (ODS 11 --- Cidades e Comunidades Sustentáveis)}, cujo propósito oficial definido pelas Nações Unidas consiste em \emph{``tornar as cidades e os assentamentos humanos inclusivos, seguros, resilientes e sustentáveis''} \cite{onubrasil2015ods11}. Como 87,4\% da população brasileira habita áreas urbanas \cite{ibge2024censo2022malha}, a sustentabilidade das cidades depende diretamente da capacidade municipal de conservar as vias públicas, garantir o funcionamento contínuo dos serviços básicos de zeladoria e instituir canais permanentes de participação cidadã no cuidado com o território \cite{ipea2019ods11}. A aderência estrutural do sistema proposto ao ODS 11 materializa-se prioritariamente nas \textbf{Metas 11.3 e 11.7}, além de dialogar diretamente com a \textbf{Meta 11.6} \cite{onubrasil2015ods11, ipea2019ods11}.

\subsection{Aderência à Meta 11.3 e ao Indicador 11.3.2 (Gestão Urbana Participativa e Integrada)}
\label{subsec:meta-11-3}

De acordo com a redação oficial da Agenda 2030 no Brasil \cite{onubrasil2015ods11}, a \textbf{Meta 11.3} estabelece o seguinte compromisso:

\begin{citacao}
\textbf{Meta 11.3} --- Até 2030, aumentar a urbanização inclusiva e sustentável, e as capacidades para o planejamento e gestão de assentamentos humanos participativos, integrados e sustentáveis, em todos os países \cite{onubrasil2015ods11}.
\end{citacao}

No quadro de monitoramento acompanhado pelo \citeonline{ipea2019ods11}, essa meta é mensurada pelo \textbf{Indicador 11.3.2}, definido como a \emph{``proporção de cidades com uma estrutura de participação direta da sociedade civil no planejamento e gestão urbana que opera de forma regular e democrática''}. O sistema proposto atende diretamente à Meta 11.3 e ao Indicador 11.3.2 ao converter cada morador cadastrado em agente ativo da gestão urbana, permitindo a abertura de ocorrências geolocalizadas por CEP com anexos de mídia e conferindo-lhe o poder de avaliar a qualidade do reparo executado. Simultaneamente, amplia a capacidade de gestão integrada das prefeituras ao conectar as demandas cidadãs ao orçamento das secretarias, às equipes de campo, à frota e ao almoxarifado.

\subsection{Aderência à Meta 11.7 e aos Indicadores 11.7.1 e 11.7.2 (Espaços Públicos Seguros e Acessíveis)}
\label{subsec:meta-11-7}

De modo convergente, a \textbf{Meta 11.7} do ODS 11 volta-se à conservação qualificada e à segurança dos logradouros de uso comum, conforme texto oficial da \citeonline{onubrasil2015ods11}:

\begin{citacao}
\textbf{Meta 11.7} --- Até 2030, proporcionar o acesso universal a espaços públicos seguros, inclusivos, acessíveis e verdes, particularmente para as mulheres e crianças, pessoas idosas e pessoas com deficiência \cite{onubrasil2015ods11}.
\end{citacao}

A Meta 11.7 associa-se ao \textbf{Indicador 11.7.1} (\emph{proporção média da área construída das cidades que é espaço público aberto para uso de todos, por sexo, idade e pessoas com deficiência}) e ao \textbf{Indicador 11.7.2} (relativo à segurança nos espaços públicos) \cite{ipea2019ods11}. Conforme revelam os dados do Censo Demográfico 2022 \cite{ibge2025entorno}, a existência física de uma via não assegura seu uso seguro e inclusivo quando o pavimento está esburacado, a iluminação pública está queimada ou os semáforos encontram-se inoperantes. Ao priorizar por gravidade o reparo de iluminação LED, pavimentação e sinalização viária, o banco de dados promove diretamente a segurança noturna e a acessibilidade nos espaços públicos.

\subsection{Articulação Complementar no ODS 11: Meta 11.6 (Resíduos e Drenagem Urbana)}
\label{subsec:meta-11-6}

Dentro do próprio ODS 11, destaca-se ainda a articulação com a \textbf{Meta 11.6} (\emph{``Até 2030, reduzir o impacto ambiental negativo per capita das cidades, inclusive prestando especial atenção à qualidade do ar, gestão de resíduos municipais e outros''} \cite{onubrasil2015ods11}), monitorada pelo \textbf{Indicador 11.6.1}. Como o sistema contempla o registro geolocalizado de pontos de descarte irregular de entulho e de obstruções na rede de drenagem pluvial \cite{snis2023drenagem, snis2023residuos}, sua operação apoia a erradicação de focos de poluição e alagamentos urbanos.

\section{Enquadramento em ODS Secundários (ODS 9 e ODS 16)}
\label{sec:ods-secundarios}

Além de sua vinculação primária ao ODS 11, o projeto estabelece sinergias diretas com dois objetivos secundários da Agenda 2030:

\begin{itemize}
    \item \textbf{ODS 9 --- Indústria, Inovação e Infraestrutura (Meta 9.1)}: cujo lema consiste em \emph{``construir infraestruturas resilientes, promover a industrialização inclusiva e sustentável e fomentar a inovação''} \cite{onubrasil2015ods9}. O sistema atua diretamente na \textbf{Meta 9.1} (\emph{``Desenvolver infraestrutura de qualidade, confiável, sustentável e resiliente, incluindo infraestrutura regional e transfronteiriça, para apoiar o desenvolvimento econômico e o bem-estar humano, com foco no acesso equitativo e a preços acessíveis para todos''} \cite{onubrasil2015ods9}), garantindo a manutenção contínua e o controle técnico de insumos, equipes e frotas que preservam a confiabilidade da infraestrutura municipal;
    \item \textbf{ODS 16 --- Paz, Justiça e Instituições Eficazes (Metas 16.6 e 16.7)}: voltado a construir instituições eficazes, responsáveis e inclusivas \cite{onubrasil2015ods16}. Ao rastrear o orçamento mensal das secretarias, os custos de mão de obra das Ordens de Serviço, os lotes de compra de fornecedores e a avaliação cidadã do atendimento, o banco de dados materializa a \textbf{Meta 16.6} (\emph{``Desenvolver instituições eficazes, responsáveis e transparentes em todos os níveis''}) e a \textbf{Meta 16.7} (\emph{``Garantir a tomada de decisão responsiva, inclusiva, participativa e representativa em todos os níveis''}) \cite{onubrasil2015ods16}, em consonância com a Lei n\textsuperscript{o}~13.460/2017 \cite{brasil2017lei13460}.
\end{itemize}

A Tabela~\ref{tab:enquadramento-ods} consolida a correspondência entre os ODS, suas metas e indicadores oficiais e as funcionalidades do sistema proposto.

\begin{table}[htb]
\ABNTEXfontereduzida
\caption{Matriz de enquadramento do problema e do sistema proposto nos Objetivos de Desenvolvimento Sustentável da Agenda 2030}
\label{tab:enquadramento-ods}
\centering
\begin{tabular}{L{2.7cm} L{1.8cm} L{4.2cm} L{5.5cm}}
\toprule
\textbf{ODS / Nível} & \textbf{Meta / Indicador} & \textbf{Foco Oficial na Agenda 2030} & \textbf{Materialização no Escopo do Sistema} \\
\midrule
\textbf{ODS 11} \newline Cidades e Comunidades Sustentáveis \newline \emph{(Principal)} & Meta 11.3 \newline Ind. 11.3.2 & Urbanização inclusiva e capacidade de planejamento e gestão participativa e integrada & Abertura de ocorrências pelo cidadão por CEP com envio de mídias, integração entre secretarias e avaliação social do serviço \\
\addlinespace
\textbf{ODS 11} \newline Cidades e Comunidades Sustentáveis \newline \emph{(Principal)} & Meta 11.7 \newline Ind. 11.7.1 \newline Ind. 11.7.2 & Acesso universal a espaços públicos seguros, inclusivos e acessíveis para grupos vulneráveis & Priorização por gravidade de reparos em iluminação pública, pavimentação, calçadas e semáforos para segurança viária e pedonal \\
\addlinespace
\textbf{ODS 11} \newline Cidades e Comunidades Sustentáveis \newline \emph{(Principal)} & Meta 11.6 \newline Ind. 11.6.1 & Redução do impacto ambiental negativo das cidades e gestão adequada de resíduos municipais & Mapeamento e remoção de pontos de descarte irregular de entulho e desobstrução de redes de drenagem pluvial \\
\midrule
\textbf{ODS 9} \newline Indústria, Inovação e Infraestrutura \newline \emph{(Secundário)} & Meta 9.1 & Desenvolvimento e conservação de infraestrutura confiável, sustentável, resiliente e de qualidade & Controle técnico de Ordens de Serviço, equipes especializadas, veículos da frota e insumos de manutenção urbana \\
\addlinespace
\textbf{ODS 16} \newline Paz, Justiça e Instituições Eficazes \newline \emph{(Secundário)} & Metas 16.6 e 16.7 & Instituições públicas eficazes, transparentes e tomada de decisão responsiva e participativa & Rastreabilidade de orçamento mensal, custos de mão de obra, lotes de fornecedores, baixa de estoque e auditoria cidadã \\
\bottomrule
\end{tabular}
\legend{Fonte: Elaborado pelos autores com base em \citeonline{onubrasil2015ods11}, \citeonline{onubrasil2015ods9}, \citeonline{onubrasil2015ods16} e \citeonline{ipea2019ods11}.}
\end{table}
